% Sampling with missing support is a mathematical impossibility example.
\clearpage
\section*{Research laboratory: a weighted estimator cannot invent missing data}
\noindent Importance sampling often rewrites an expectation as
\[
\mathbb E_p[f(X)]
=\sum_{x:q(x)>0}q(x)\frac{p(x)}{q(x)}f(x).
\]
That equality requires \(p(x)f(x)=0\) wherever \(q(x)=0\);
for general expectations the proposal must assign positive
probability to every state having positive target probability.
Otherwise parts of the target integral are omitted and no
amount of renormalizing the observed weights can restore them.

For a concrete two-state example, let the target be
\(p(0)=p(1)=1/2\), and let \(f(0)=0,\ f(1)=1\).
The true expectation is \(1/2\).
A proposal \(q(0)=1,\ q(1)=0\) will \emph{never}
observe state 1. Its sampled weighted contributions are all zero;
a weighted average built from those samples remains zero
regardless of the number of repetitions.
Clipping weights or running more trials under the same proposal
does not solve the support failure.

When the support condition holds, other risks remain:
a proposal assigning very small probability to important events
can produce enormous weights and variance, while reused samples,
adaptive selection and correlated trajectories complicate the
interpretation of confidence intervals. Reproducible random seeds
do not guarantee that observations constitute independent
statistical units.

\paragraph{Reconstruction exercise.}
Implement the two-state example with an explicit support check.
Change the proposal to \(q(1)=0.01\), then compare the sample
average, importance-weighted estimate and empirical variance over
many independently seeded experiments. Report how often rare-event
coverage actually occurs, and separate sampling error from bias.
Only after those checks should the technique be transported to
image sampling, geometry reconstruction or research benchmarks.
